
\textbf{Proposition 2.2.} \uline{Avec les notations et les hypothèses de} 2.1, \uline{soit} $f:P\to S$ \uline{le morphisme structural de} $P$. \uline{Alors l'homomorphisme} :
\[
f^*:K_\bullet(S)\longrightarrow K_\bullet(P)
\]
\uline{est surjectif}, \uline{et son noyau est le sous-groupe de} $K_\bullet(S)$ \uline{engendré par les sous-groupes} $(\eta(m)-1)K_\bullet(S)$ \uline{pour} $m\in M$, \uline{ou}, \uline{ce qui revient au même}, \uline{pour} $m$ \uline{parcourant un système de générateurs de} $M$.

On suppose d'abord $r=1$, i.e.\ $T=G_{m\,S}$. Alors $M=\mathbb Z$, et on a pour tout $m\in M$ :
\[
\underline{A}_m\simeq \underline{A}_1^{\otimes m}.
\]
Par suite, si on pose $\underline{A}_1=\underline{L}$, on a $\underline{A}=\bigoplus_{n\in\mathbb Z}\underline{L}^{\otimes n}$, ce qui montre que $P$ s'identifie au complémentaire de la section nulle dans le fibré $X=\mathbb V(\underline{L})$. Considérant $S$ comme un sous-schéma fermé de $X$ par la section nulle, on a d'après 1.1 la suite exacte :
\[
K_\bullet(S)\xrightarrow{g_*}K_\bullet(X)\longrightarrow K_\bullet(P)\longrightarrow 0
\]
où $g:S\to X$ est la section nulle de $X$. Par ailleurs, $f^*$ se factorise en \(K_\bullet(S)\xrightarrow{f'^*}K_\bullet(X)\longrightarrow K_\bullet(P)\) où $f'$ est le morphisme structural de $X$. D'après 1.6, $f'^*$ est un isomorphisme, inverse de $g^*$. On a donc la suite exacte :
\[
K_\bullet(S)\xrightarrow{g^*\circ g_*}K_\bullet(S)\xrightarrow{f^*}K_\bullet(P)\longrightarrow 0,
\]
d'où l'on tire d'après 1.8 :
\[
K_\bullet(P)\simeq K_\bullet(S)/\lambda_{-1}(L)K_\bullet(S)=K_\bullet(S)/(1-L)K_\bullet(S)
\]
où $L$ est la classe dans $K^{\bullet}(S)$ de $\underline{L}$. Comme pour tout $n\in\mathbb Z$, $L^n-1$ est multiple de $1-L$, ceci prouve 2.2 dans le cas où $P$ est de rang $1$.

